\section{Evolución del sistema de evaluación: de MBPP/HumanEval a SWE-Bench}

\subsection{El valor y la caducidad de las evaluaciones tempranas de código}
MBPP y HumanEval impulsaron en su momento la iteración rápida de los modelos de código, porque cumplían entonces tres condiciones clave: dificultad adecuada, calificación automática y una filtración aún no generalizada. Sutton señala directamente que hoy este tipo de tareas se parece más al ``MNIST de la evaluación de código'': todavía sirve para comparaciones básicas, pero difícilmente representa la capacidad real de ingeniería de software.

\begin{knowledgebox}{Interpretación de Pass@K y sus implicaciones de despliegue}
Pass@K mide la probabilidad de que ``al muestrear K veces, al menos una pase las pruebas''. Tiene valor para evaluar la diversidad, pero no siempre equivale a la experiencia del producto:
\begin{itemize}
    \item la finalización única en el IDE se acerca más a Pass@1;
    \item un agente con una etapa de verificación puede acercarse más a Pass@K;
    \item por lo tanto, si la métrica es válida depende de cómo se use realmente el sistema.
\end{itemize}
\end{knowledgebox}

\begin{warningbox}{Un desajuste de métricas provoca un rumbo equivocado}
Si el escenario del producto solo permite una salida, pero el equipo optimiza principalmente Pass@K, puede terminar entrenando el sistema para que ``el conjunto de muestras contenga la respuesta correcta'', en lugar de que ``la primera respuesta sea utilizable''. Este sesgo se vuelve muy evidente en la etapa de lanzamiento.
\end{warningbox}

\subsection{El avance de SWE-Bench}
SWE-Bench hizo avanzar la tarea de ``escribir funciones cortas'' a ``corregir issues reales en un repositorio de código real'': la entrada es la descripción del issue y el contexto del repositorio, la salida es un patch, y la verificación depende del conjunto de pruebas ya existente. En comparación con las evaluaciones tempranas, elevó de forma notable el realismo y obligó a los investigadores a enfrentar problemas de sistema como la navegación de código, la recuperación de contexto y la edición de múltiples archivos.

\begin{figure}[H]
\centering
\includegraphics[width=0.88\textwidth]{frame-02-swebench.jpg}
\caption{Captura de la clase donde se explica la forma de la tarea de SWE-Bench y su impacto \protect\footnotemark}
\end{figure}
\footnotetext{Intervalo de tiempo en el video: 00:14:55--00:15:40.}

\begin{importantbox}{Lo que SWE-Bench realmente cambió}
Lo que cambió no fue solo la puntuación, sino el paradigma de investigación y desarrollo. Los equipos comenzaron a invertir de manera sistemática en:
\begin{itemize}
    \item recuperación y navegación a nivel de repositorio de código;
    \item corrección impulsada por retroalimentación de ejecución;
    \item gestión de trayectorias de múltiples turnos;
    \item un ciclo cerrado de generación y verificación de patches.
\end{itemize}
\end{importantbox}

\subsection{Las limitaciones de SWE-Bench son igual de importantes}
Sutton lo dice con franqueza: el valor y las limitaciones de SWE-Bench coexisten. Las fuentes de datos se concentran en un número reducido de proyectos y ecosistemas de lenguajes específicos; el estilo de las descripciones de los issues no coincide del todo con la interacción real entre humano y máquina dentro del IDE; los commits históricos tienen el ruido de ``un commit que mezcla varios cambios''; el conjunto de pruebas no es un verifier perfecto, y pueden aparecer parches que ``sobreajustan las pruebas pero son semánticamente incorrectos''.

\begin{table}[H]
\centering
\begin{tabular}{p{3.4cm}p{5.0cm}p{5.0cm}}
\toprule
\textbf{Dimensión} & \textbf{Contribución de SWE-Bench} & \textbf{Sesgo potencial} \\
\midrule
Realismo de la tarea & Proviene de la cadena issue/patch de proyectos reales & Distribución de proyectos limitada, cobertura de dominios incompleta \\
Forma de verificación & Puede ejecutar pruebas automáticamente y formar una evaluación de ciclo cerrado & Las pruebas no representan completamente la corrección semántica \\
Entrada de lenguaje & Conserva el estilo de redacción real de los issues y las pistas de contexto & Existe una diferencia de distribución respecto a las instrucciones breves dentro del IDE \\
Valor de entrenamiento & Puede impulsar el diseño de retrieval, editing y execution & Los commits históricos pueden entrelazar múltiples objetivos de cambio \\
\bottomrule
\end{tabular}
\caption{El ``fuerte realismo'' y el ``fuerte sesgo'' de SWE-Bench coexisten}
\end{table}

\subsection{La vida útil de las evaluaciones y su renovación continua}
La clase plantea un juicio muy importante: \textbf{toda evaluación tiene una shelf life}. Al principio es difícil y no se ha filtrado, pero a los pocos años suele volverse fácil de resolver y de memorizar. El sistema de evaluación debe actualizarse de manera continua, o el sistema terminará optimizándose eficientemente sobre un objetivo caduco.

\begin{importantbox}{Una frase metodológica que se puede poner en práctica}
\textit{``If your numbers aren't good, you're hill climbing on the wrong thing.''}  
Traducido a una acción de ingeniería, esto significa: primero verificar si el objetivo de evaluación es real y si todavía tiene poder de discriminación, y solo después ajustar los parámetros del modelo o del harness.
\end{importantbox}

\subsection{Resumen de la sección}
De MBPP/HumanEval a SWE-Bench, la industria no se limitó a ``cambiar de tabla de clasificación'', sino que llevó la inteligencia de código del nivel de función al nivel de repositorio. Al mismo tiempo, la evaluación misma debe auditarse de manera continua, o una puntuación alta no equivaldrá a una alta utilidad práctica.

